\medskip
\noindent\emph{Additional audit of the residual test
$k=\mathbb Q^{\mathrm{solv}}$.}
Neither of the elementary nongeometric constructions in F.49 settles
this test. In fact, with $G=G_{\mathbb Q}$ and
$H=G_{\mathbb Q^{\mathrm{solv}}}$, one has
\[
                         C_G(H)=1.                         \tag{F44.C1}
\]
We give the proof, verifying the precise scope of the recovered
\texttt{CYCLOTOMIC\_KERNEL} row B2 instead of importing a stronger
unqualified assertion about arbitrary closed subgroups.

The following argument works for every Hilbertian field $F$ and
every nontrivial closed normal subgroup $N\triangleleft G_F$.
Two primary results suffice. Jarden's Proposition~2.5 says that if
$E/F$ is algebraic Galois and $E$ is not separably closed, then
$G_E$ is almost locally free; in particular it contains a free
profinite subgroup of rank two
\cite{F44Jarden2011}. Haran--Jarden, Corollary~2.5, says that
$F^{\mathrm{sep}}$ cannot be the compositum of two proper subfields
which are Galois over $F$ \cite{F44HaranJarden1991}. Both statements
allow infinite Galois extensions.

There is no nontrivial closed abelian normal subgroup $A$ of $G_F$:
its fixed field is Galois over $F$ and is not separably closed, so
the first result would place a nonabelian free profinite rank-two
group inside $A$. Put $C=C_{G_F}(N)$. It is closed and normal,
and $N\cap C=Z(N)$ is a closed abelian normal subgroup of $G_F$;
hence $N\cap C=1$. If $C\ne1$, the two proper fixed fields of
$N$ and $C$ would be Galois over $F$, and their compositum would
have absolute Galois group $N\cap C=1$. This contradicts the
second result. Thus $C_{G_F}(N)=1$.

To apply it to $F=\mathbb Q$, we check that
$\mathbb Q^{\mathrm{solv}}\ne\overline{\mathbb Q}$.
The polynomial $P=X^5-X-1$ is irreducible modulo $3$: the attached
exact program divides it by all twelve monic polynomials of degree
one or two and finds no divisor. A reducible quintic must have a
factor of one of these degrees. Modulo $2$ its squarefree
factorization is
\[
           (X^2+X+1)(X^3+X^2+1),
\]
with both factors irreducible. Its rational Galois group is therefore
transitive of degree five and contains a permutation of cycle type
$(2)(3)$, whose cube is a transposition. A transitive group of prime
degree is primitive. The graph whose edges are the conjugates of
that transposition has components forming blocks, and is therefore
connected. Its edge transpositions generate $S_5$, so the Galois
group is $S_5$. The program also computes its derived-subgroup
orders $120,60,60$, certifying nonsolvability. Its splitting field
cannot lie in the maximal solvable extension. Thus $H\ne1$, and
Hilbert irreducibility for $\mathbb Q$ gives (F44.C1).
In particular $Z(H)=1$, so F.49 supplies no nonidentity central
comparison twist over this field.

The semilinear field-automorphism construction also does not apply.
Since $k/\mathbb Q$ is Galois,
$\operatorname{Aut}_{\mathbb Q}(k)=G/H$. For an automorphism
$\sigma$ choose a lift $g\in G$. Its outer action on $H$ is
conjugation by $g$. If this action were inner, say equal to
conjugation by $h\in H$, then
\[
                       h^{-1}g\in C_G(H)=1.
\]
Hence $g=h$ and $\sigma=1$. A nonidentity coefficient-semilinear
automorphism of $\mathbb A^1_k$ consequently cannot be supplied with
the relative comparison used in F.49: such a comparison, restricted
along the rational point $0$, would make its base automorphism of
$BH$ two-isomorphic to the identity and thus its action on $H$
inner. This excludes that construction; it makes no assertion that
all pinned topoi morphisms over $k$ have been classified.

For orientation on the positive arguments, $k^*$ is divisible.
An element of $k$ lies in a finite solvable Galois extension of
$\mathbb Q$. Adjoin the $n$th roots of unity and all $n$th roots
of its conjugates. The new Galois extension is abelian over the
intermediate field and solvable over $\mathbb Q$, so the desired
roots remain in $k$. Thus $k^*/\mu_\infty$ is a nonzero divisible group (the class of
$2$ is nonzero), so it cannot be free abelian. The corresponding
input of F.50 cannot be used here. Also $k$ cannot embed over $\mathbb Q$
in any finitely generated extension of $\mathbb Q^{\mathrm{ab}}$:
being algebraic over $\mathbb Q$, its image would lie in the finite
constant extension, whose uniform positive non-torsion height bound
contradicts all the elements $2^{1/n}$. These are precise limitations
of the arguments of this return. The remaining pinned Hom-question
over $\mathbb Q^{\mathrm{solv}}$ stays unresolved here, without a
claim that the corresponding question is open in the entire
literature.

The finite certificate is
\texttt{code/F44\_solvable\_certificates.py}; its exact output is
included with the deliverable.
